\documentclass[8pt, a4paper]{article}

% ── Packages ─────────────────────────────────────────────────────────────────
\usepackage[a4paper, top=2.5cm, bottom=2.5cm, left=2.8cm, right=2.8cm]{geometry}
\usepackage[T1]{fontenc}
\usepackage[utf8]{inputenc}
\usepackage{lmodern}
\usepackage{microtype}
\usepackage{parskip}          % space between paragraphs, no indent
\usepackage{titlesec}         % section heading styles
\usepackage{enumitem}         % list customisation
\usepackage{booktabs}         % professional tables
\usepackage{tabularx}         % auto-width table columns
\usepackage{array}            % extra column formatting
\usepackage{xcolor}           % colours
\usepackage{mdframed}         % framed boxes
\usepackage{fancyhdr}         % header / footer
\usepackage{listings}         % code listings
\usepackage{hyperref}         % clickable links (load last)

% ── Colours ───────────────────────────────────────────────────────────────────
\definecolor{accentblue}{RGB}{30, 90, 160}
\definecolor{rulegray}{RGB}{190, 195, 205}
\definecolor{boxbg}{RGB}{237, 242, 250}
\definecolor{codebg}{RGB}{245, 247, 250}
\definecolor{codeframe}{RGB}{180, 190, 210}

% ── Hyperref setup ────────────────────────────────────────────────────────────
\hypersetup{
  colorlinks = true,
  linkcolor  = accentblue,
  urlcolor   = accentblue,
}

% ── Section heading styles ────────────────────────────────────────────────────
\titleformat{\section}
  {\large\bfseries\color{accentblue}}
  {\thesection.}{0.6em}{}
  [\vspace{-4pt}\textcolor{rulegray}{\rule{\linewidth}{0.6pt}}\vspace{2pt}]

\titlespacing{\section}{0pt}{18pt}{6pt}

% ── Concept box ───────────────────────────────────────────────────────────────
\newmdenv[
  backgroundcolor   = boxbg,
  linecolor         = accentblue,
  linewidth         = 1.2pt,
  leftline          = true,
  rightline         = false,
  topline           = false,
  bottomline        = false,
  innerleftmargin   = 10pt,
  innertopmargin    = 6pt,
  innerbottommargin = 6pt,
  skipabove         = 6pt,
  skipbelow         = 6pt,
]{conceptbox}

% ── Code listings ─────────────────────────────────────────────────────────────
\lstset{
  basicstyle       = \ttfamily\footnotesize,
  backgroundcolor  = \color{codebg},
  frame            = single,
  rulecolor        = \color{codeframe},
  framesep         = 4pt,
  xleftmargin      = 8pt,
  xrightmargin     = 8pt,
  breaklines       = true,
  breakatwhitespace= true,
  showstringspaces = false,
  tabsize          = 4,
}

% ── List styling ──────────────────────────────────────────────────────────────
\setlist[itemize]{topsep=3pt, itemsep=2pt, parsep=0pt, leftmargin=1.6em}
\setlist[enumerate]{topsep=3pt, itemsep=2pt, parsep=0pt, leftmargin=1.8em}

% ── Header / Footer ───────────────────────────────────────────────────────────
\pagestyle{fancy}
\fancyhf{}
\renewcommand{\headrulewidth}{0.4pt}
\fancyhead[L]{\small\color{accentblue}\textbf{ELEC3609 / ELEC9609}}
\fancyhead[R]{\small\color{gray}Week 4 --- Front-End Development \& the Browser Ecosystem}
\fancyfoot[C]{\small\thepage}

% =============================================================================
\begin{document}

% ── Title Block ───────────────────────────────────────────────────────────────
\begin{center}
  {\LARGE\bfseries\color{accentblue} ELEC3609 / ELEC9609: Internet Software Platforms}\\[6pt]
  {\large Week 4 Study Guide: Front-End Development \& the Browser Ecosystem}
\end{center}

\vspace{12pt}
\textcolor{rulegray}{\rule{\linewidth}{1pt}}
\vspace{6pt}

% =============================================================================
\section{Frontend Fundamentals \& User Experience (UX vs UI)}

The front-end is the presentation layer of a web application. It executes directly in the client browser,
serving as the bridge between the user and backend services (data persistence, business logic, computations).

\begin{conceptbox}
\textbf{Core Usability Philosophy}
\begin{itemize}
  \item \textbf{Usability Axiom:} ``Mobile is a magnifying glass for your usability problems.''
  \item \textbf{Engineering vs.\ Art:} Building effective user interfaces is NOT art---it is an engineering
        discipline based on structure, sensible defaults, and consistency.
  \item \textbf{Steve Jobs on Design:} ``Design is not just what it looks like and feels like. Design is how it works.''
\end{itemize}
\end{conceptbox}

\begin{conceptbox}
\textbf{UX (User Experience) vs.\ UI (User Interface)}

UX and UI represent two distinct aspects of frontend design:

\smallskip
\begin{tabularx}{\linewidth}{l X}
  \toprule
  \textbf{Role} & \textbf{Key Responsibilities \& Focus Areas} \\
  \midrule
  \textbf{UX Designer} (Interaction Design)
    & Focuses on HOW IT WORKS \& user interaction flow. Charts user pathways \& plans information architecture.
      Creates structural wireframes \& interactive prototypes; conducts user research. \\[4pt]
  \textbf{UI Designer} (Visual Design)
    & Focuses on HOW IT LOOKS \& visual aesthetics. Selects color palettes, typography, \& styling.
      Produces high-fidelity mockups, graphics, and pixel-perfect page layouts. \\[4pt]
  \textbf{Technical Artist}
    & Bridges the gap between visual artwork and frontend code execution/performance. \\
  \bottomrule
\end{tabularx}
\end{conceptbox}

\begin{conceptbox}
\textbf{Golden Rules for Great UX}
\begin{enumerate}
  \item \textbf{No Instructions Needed:} The interface should be intuitive enough that users do not need a user manual.
  \item \textbf{Simple Common Tasks:} Keep the user path for frequent actions as short as possible.
  \item \textbf{Every Click Matters:} Minimize unnecessary steps and friction.
  \item \textbf{User Happiness:} Design predictable, responsive interactions that reduce user frustration.
\end{enumerate}
\end{conceptbox}

% =============================================================================
\section{Webpage Design Pipeline: Wireframes, Mockups \& Prototypes}

The webpage design workflow moves sequentially from structural outline to high-fidelity interactive behavior.

\begin{conceptbox}
\textbf{The Three Design Stages}
\begin{enumerate}
  \item \textbf{Wireframe (Structural Blueprint):}
    \begin{itemize}
      \item \textit{Scope:} Structural layout and component positioning.
      \item \textit{Purpose:} Outline what information and UI components will be displayed on each page.
      \item \textit{Characteristics:} Low fidelity, grayscale/skeletal representation; no visual styling.
    \end{itemize}
  \item \textbf{Mockup (Visual Specification):}
    \begin{itemize}
      \item \textit{Scope:} Visual appearance and aesthetics.
      \item \textit{Purpose:} Show exact colors, typography, spacing, icons, and visual branding.
      \item \textit{Characteristics:} High-fidelity static graphical snapshot of the final page design.
    \end{itemize}
  \item \textbf{Prototype (Functional Experience):}
    \begin{itemize}
      \item \textit{Scope:} Interaction, motion, and behavior.
      \item \textit{Purpose:} Demonstrate how the page works and responds to user inputs.
      \item \textit{Characteristics:} Interactive model (can be low or high fidelity) allowing users to
            click through and evaluate real workflow behavior.
    \end{itemize}
\end{enumerate}
\end{conceptbox}

\begin{conceptbox}
\textbf{Design Deliverables Summary}

\smallskip
\begin{tabularx}{\linewidth}{l l X}
  \toprule
  \textbf{Design Stage} & \textbf{Primary Focus} & \textbf{Key Question Answered} \\
  \midrule
  Wireframe  & Information Architecture \& Layout & ``Where do components go?'' \\
  Mockup     & Visual Style \& Graphic Aesthetic  & ``What will it look like?'' \\
  Prototype  & Interaction \& Behavioral Flow     & ``How does the user interact?'' \\
  \bottomrule
\end{tabularx}
\end{conceptbox}

% =============================================================================
\section{The Browser Ecosystem, Interoperability \& Cross-Compatibility}

\begin{conceptbox}
\textbf{Browser Fragmentation Challenges}

Web applications execute in highly heterogeneous runtime environments:
\begin{itemize}
  \item \textbf{Browser Engines:} Chrome (Blink), Firefox (Gecko), Safari (WebKit), Edge (Blink).
  \item \textbf{Versioning:} Legacy versions vs.\ evergreen auto-updating browsers.
  \item \textbf{Operating Systems:} Windows, macOS, Linux, iOS, Android (each rendering differently).
  \item \textbf{Hardware Diversity:} Screen resolutions, high-DPI viewports, touch vs.\ mouse inputs.
\end{itemize}
\end{conceptbox}

\begin{conceptbox}
\textbf{Interoperability Strategy}
\begin{itemize}
  \item \textbf{Rule of Support:} You do NOT need to support all browsers, versions, or operating systems.
  \item \textbf{Strategic Targeting:} At project inception, define target supported browsers based strictly
        on target audience metrics and demographic analytics.
  \item \textbf{Compatibility Reference:} \texttt{caniuse.com} provides up-to-date compatibility tables for
        HTML5, CSS, and JS APIs across browser versions.
\end{itemize}
\end{conceptbox}

\begin{conceptbox}
\textbf{Testing Methodologies}
\begin{enumerate}
  \item \textbf{Local Machine Testing:} Native execution in locally installed browser instances.
  \item \textbf{Virtual Machine (VM) Emulation:} Running target OS environments inside VMs.
        \textit{Note: Virtualization can introduce its own unique rendering/performance bugs.}
  \item \textbf{Professional Cloud Platforms:} Automated cross-browser testing services running on real
        physical hardware/devices (e.g., BrowserStack, Sauce Labs).
\end{enumerate}
\end{conceptbox}

% =============================================================================
\section{Modern Design Practices \& Web Components}

\begin{conceptbox}
\textbf{Design Best Practices}
\begin{enumerate}
  \item \textbf{Responsive Web Design (RWD):} Fluid page layouts that dynamically adapt to any screen size
        or orientation using media queries, flexible grids, and responsive media.
  \item \textbf{Inclusive Design:} Creating software accessible to everyone without exclusion. Considers
        accessibility across: Disability (visual, auditory, motor), age, language, ethnicity, race,
        religion, gender identity, and sexual orientation.
\end{enumerate}
\end{conceptbox}

\begin{conceptbox}
\textbf{Web Components Architecture}

Web Components are a suite of native browser technologies allowing developers to create custom, reusable
HTML tags (\texttt{<user-card>}, \texttt{<custom-button>}) containing HTML, CSS, and JS.

\medskip
\textbf{The 4 Pillars of Web Components}
\begin{enumerate}
  \item \textbf{Reactivity:} Automatically re-rendering or updating UI state when properties/attributes change.
  \item \textbf{Encapsulation:} Keeping styles and markup scoped internally so they don't leak out (Shadow DOM).
  \item \textbf{Modularity:} Packaging complex UI widgets into standalone, self-contained units.
  \item \textbf{Reusability:} Using the custom element repeatedly across pages or projects like native HTML tags.
\end{enumerate}
\end{conceptbox}

% =============================================================================
\section{HTML5 Markup, Document Structure \& the DOM}

\begin{conceptbox}
\textbf{HTML Specification Timeline}
\begin{itemize}
  \item \textbf{1997:} HTML 3.2
  \item \textbf{1999:} HTML 4.01
  \item \textbf{2007:} HTML5 Draft published
  \item \textbf{2014:} HTML5 published as official W3C Recommendation (introduced \texttt{<audio>}, \texttt{<video>},
        semantic elements, and input types: \texttt{color}, \texttt{date}, \texttt{email}, \texttt{time}).
\end{itemize}
\end{conceptbox}

\begin{conceptbox}
\textbf{The DOCTYPE Declaration}
\begin{itemize}
  \item \textbf{Purpose:} NOT an HTML tag. It is a directive informing the browser parser which standard or
        HTML version is being used.
  \item \textbf{Placement:} Must be the absolute FIRST line in the HTML document.
  \item \textbf{HTML5 DOCTYPE:} \texttt{<!DOCTYPE html>} --- short, case-insensitive, standards-compliant,
        and masquerades safely for older browsers.
  \item \textbf{Legacy HTML4 DOCTYPE (Complex DTD):}\\
        \texttt{<!DOCTYPE HTML PUBLIC ``-//W3C//DTD HTML 4.01//EN'' ``http://www.w3.org/TR/html4/strict.dtd''>}
\end{itemize}
\end{conceptbox}

\begin{conceptbox}
\textbf{Minimal HTML5 Document Template}

\begin{lstlisting}[language=HTML]
<!DOCTYPE html>
<html lang="en">
<head>
    <meta charset="UTF-8">
    <title>Page Title</title>
    <link rel="stylesheet" href="styles.css">
</head>
<body>
    <h1>Main Heading</h1>
    <p>Paragraph content goes here.</p>

    <!-- Scripts placed at end of body to avoid blocking page render -->
    <script src="app.js"></script>
</body>
</html>
\end{lstlisting}
\end{conceptbox}

\begin{conceptbox}
\textbf{Structural HTML Tags Breakdown}
\begin{itemize}
  \item \textbf{\texttt{<html>}:} Root element wrapping the entire document. Defines language attribute (\texttt{lang="en"}).
  \item \textbf{\texttt{<head>}:} Stores document metadata, encoding, title, and external CSS style links.
    \begin{itemize}
      \item \textit{Note:} Script tags in \texttt{<head>} block page parsing/rendering while downloading.
      \item \textit{Mandatory Tag:} \texttt{<title>} is required inside \texttt{<head>}.
    \end{itemize}
  \item \textbf{\texttt{<body>}:} Contains all visible UI content rendered in the browser viewport.
    \begin{itemize}
      \item \textit{Best Practice:} Place \texttt{<script>} tags immediately before closing \texttt{</body>} tag to
            ensure DOM elements are parsed before scripts execute.
      \item \textit{Empty Elements:} Self-closing tags without closing tags (e.g., \texttt{<br>}, \texttt{<img>}).
    \end{itemize}
\end{itemize}

\smallskip
\textbf{HTML Attribute Best Practices}
\begin{itemize}
  \item Use lowercase for all element names and attributes (\texttt{src}, \texttt{href}, \texttt{title}, \texttt{style}).
  \item Always enclose attribute values in quotes (\texttt{href="https://sydney.edu.au"}).
\end{itemize}
\end{conceptbox}

\begin{conceptbox}
\textbf{Document Object Model (DOM)}
\begin{itemize}
  \item \textbf{Concept:} An in-memory object-oriented tree structure representing the HTML document.
  \item \textbf{DOM Manipulation:} JavaScript interacts with the DOM to dynamically alter page content,
        styles, and attributes.
\end{itemize}

\smallskip
\textbf{Common DOM Selection Methods}

\smallskip
\begin{tabularx}{\linewidth}{l X}
  \toprule
  \textbf{DOM Method} & \textbf{Description / Return Value} \\
  \midrule
  \texttt{document.getElementById('id')}        & Returns single matching element by ID. \\
  \texttt{document.getElementsByTagName('tag')} & Returns HTMLCollection of elements by tag. \\
  \texttt{document.getElementsByClassName('c')} & Returns HTMLCollection of elements by class. \\
  \texttt{document.querySelector('selector')}   & Returns first element matching CSS selector. \\
  \texttt{document.querySelectorAll('sel')}      & Returns NodeList of all matching elements. \\
  \bottomrule
\end{tabularx}
\end{conceptbox}

\begin{conceptbox}
\textbf{Web Component Custom Element Lifecycle Callbacks}

Custom HTML elements hook into browser DOM events via lifecycle callbacks:
\begin{itemize}
  \item \textbf{\texttt{connectedCallback()}:} Invoked when the custom element is attached to the document DOM.
  \item \textbf{\texttt{disconnectedCallback()}:} Invoked when the custom element is detached/removed from the DOM.
\end{itemize}
\end{conceptbox}

% =============================================================================
\section{CSS Architecture, Selectors, Responsive Design \& Frameworks}

\begin{conceptbox}
\textbf{Cascading Style Sheets (CSS)}
\begin{itemize}
  \item \textbf{Role:} Defines presentation and styling of HTML markup, establishing separation of concerns
        between document content (HTML) and visual appearance (CSS).
  \item \textbf{Syntax:} \texttt{selector \{ property: value; \}}
\end{itemize}
\end{conceptbox}

\begin{conceptbox}
\textbf{Primary Selectors \& Specificity}
\begin{enumerate}
  \item \textbf{Element Selector:} Targets all elements of a specific HTML tag.
\begin{lstlisting}[language=CSS]
body { background-color: #f4f4f4; font-family: Arial, sans-serif; }
\end{lstlisting}
  \item \textbf{Class Selector (\texttt{.className}):} Targets elements with a specific \texttt{class} attribute.
        Reusable; can be applied to multiple elements. Elements can have multiple classes
        (\texttt{<div class="card shadow primary">}).
\begin{lstlisting}[language=CSS]
.dog { color: blue; font-weight: bold; }
\end{lstlisting}
  \item \textbf{ID Selector (\texttt{\#idName}):} Targets a single element with a specific unique \texttt{id}.
        Must be unique; an ID can only be assigned to one element per document.
\begin{lstlisting}[language=CSS]
#dog-name { color: red; text-decoration: underline; }
\end{lstlisting}
\end{enumerate}

\smallskip
\textbf{Selector Specificity Order (Highest to Lowest Priority):}\\
Inline Styles (\texttt{style=""}) $>$ ID (\texttt{\#id}) $>$ Class (\texttt{.class}) $>$ Element (\texttt{tag})
\end{conceptbox}

\begin{conceptbox}
\textbf{Responsive Design \& Media Queries}

Media Queries (\texttt{@media}): Apply CSS rules conditionally depending on viewport dimensions, device
orientation, or screen resolution.

\begin{lstlisting}[language=CSS]
/* Default styles for mobile viewports */
.container { width: 100%; }

/* Desktop viewport override */
@media (min-width: 768px) {
    .container { width: 750px; }
}
\end{lstlisting}
\end{conceptbox}

\begin{conceptbox}
\textbf{CSS Frameworks}
\begin{itemize}
  \item \textbf{Philosophy:} ``The best CSS is written by someone else.'' Pre-built, battle-tested CSS frameworks
        accelerate development and ensure cross-browser layout consistency.
  \item \textbf{Popular Frameworks:} Bootstrap, Foundation, Materialize.
  \item \textbf{Framework Selection Criteria:} Visual aesthetics, grid layout engine, community support,
        and cross-browser compatibility.
\end{itemize}
\end{conceptbox}

% =============================================================================
\section{JavaScript Ecosystem: TypeScript, SPAs \& AJAX}

\begin{conceptbox}
\textbf{JavaScript Fundamentals}
\begin{itemize}
  \item \textbf{Characteristics:} High-level, dynamic, untyped (dynamically typed), interpreted programming language.
  \item \textbf{Function:} Manipulates the HTML DOM dynamically to create interactive client-side web apps.
\end{itemize}
\end{conceptbox}

\begin{conceptbox}
\textbf{Historical JavaScript Issues}
\begin{itemize}
  \item \textbf{Type Coercion:} Weird implicit conversions (\texttt{5 == "5"} evaluates to true).
  \item \textbf{Context Loss:} Complex and unintuitive behavior of the \texttt{this} keyword.
  \item \textbf{Weak Tooling:} Absence of static type checking leads to runtime type errors.
  \item \textbf{Lack of Native Modules:} Early JS lacked standard module importing (\texttt{import}/\texttt{export}).
\end{itemize}
\end{conceptbox}

\begin{conceptbox}
\textbf{TypeScript}
\begin{itemize}
  \item \textbf{Definition:} Open-source language developed by Microsoft; a strict superset of JavaScript.
  \item \textbf{Transpilation:} TypeScript code is transpiled back into standard JavaScript for browser execution.
  \item \textbf{Key Additions:} Optional static typing, interface contracts, classes, and ES6 module namespaces.
  \item \textbf{Target:} Designed specifically for engineering large-scale enterprise web applications.
\end{itemize}
\end{conceptbox}

\begin{conceptbox}
\textbf{jQuery}
\begin{itemize}
  \item \textbf{Legacy Role:} Historical DOM abstraction library (``Swiss army knife'') that simplified DOM
        manipulation and cross-browser compatibility.
  \item \textbf{Modern Decline:} Declining popularity because modern browsers now natively support clean DOM
        APIs (\texttt{document.querySelector}) and modern JavaScript syntax (ES6+).
\end{itemize}
\end{conceptbox}

\begin{conceptbox}
\textbf{Single Page Applications (SPAs)}

Loads a single HTML document; dynamically updates user views client-side using JavaScript without performing
full page browser reloads.

\smallskip
\begin{tabularx}{\linewidth}{X X}
  \toprule
  \textbf{SPA Advantages (Pros)} & \textbf{SPA Disadvantages (Cons)} \\
  \midrule
  Fluid, fast, app-like user interface           & Higher initial codebase complexity \& overhead \\
  Strict separation of frontend UI \& backend API & Pushes processing workload into client browser \\
  Easy deployment to static CDNs (fast load)     & Poor performance on low-spec/mobile hardware \\
  Simplified dynamic data binding to views       & SEO indexing challenges \& framework churn \\
  \bottomrule
\end{tabularx}
\end{conceptbox}

\begin{conceptbox}
\textbf{AJAX (Asynchronous JavaScript and XML)}

Fetches data asynchronously from backend servers without requiring a full page refresh.

\smallskip
\textbf{API Evolution:}
\begin{enumerate}
  \item \textbf{\texttt{XMLHttpRequest (XHR)}:} Original low-level browser API (clunky syntax).
  \item \textbf{\texttt{jQuery \$.ajax()}:} Simplified library wrapper around XHR.
  \item \textbf{\texttt{Fetch API}:} Modern, promise-based native browser API
        (\texttt{fetch('/api/data').then(...)}).
\end{enumerate}
\end{conceptbox}

\begin{conceptbox}
\textbf{Popular Frontend Frameworks}

React (Meta), Angular (Google), Vue.js, Ember.js, Backbone.js.
\end{conceptbox}

% =============================================================================
\section{Frontend Tooling \& Build Ecosystem (Bower, Task Runners, Webpack)}

\begin{conceptbox}
\textbf{Package Management (Bower)}
\begin{itemize}
  \item \textbf{Function:} Front-end package manager dedicated to HTML, CSS, and JS client-side dependencies.
  \item \textbf{Runtime:} Node.js application installed via \texttt{npm}.
  \item \textbf{Manifest:} Tracks client-side project dependencies inside \texttt{bower.json}.
\end{itemize}
\end{conceptbox}

\begin{conceptbox}
\textbf{Task Runners (Grunt \& Gulp)}
\begin{itemize}
  \item \textbf{Function:} Automated task runners (frontend equivalent to a Unix \texttt{Makefile}).
  \item \textbf{Primary Use Cases:}
    \begin{enumerate}
      \item \textbf{Minification:} Compressing HTML, CSS, JS, and image assets to reduce network transfer sizes.
      \item \textbf{Concatenation:} Combining multiple source files into unified bundles.
      \item \textbf{Linting:} Code style and syntax error checking.
      \item \textbf{Testing:} Automating client-side test execution.
    \end{enumerate}
\end{itemize}
\end{conceptbox}

\begin{conceptbox}
\textbf{Module Bundlers (Webpack)}
\begin{itemize}
  \item \textbf{Function:} Static module bundler for modern JavaScript applications.
  \item \textbf{Behavior:} Recursively builds a dependency graph of all assets (JS, CSS, images) and packages
        them into optimized, production-ready bundle files containing only necessary code.
\end{itemize}
\end{conceptbox}

\vspace{14pt}
\textcolor{rulegray}{\rule{\linewidth}{1pt}}

\end{document}
